\chapter{Mô hình hóa bài toán tìm kiếm (Problem Formulation)}

\newpage

\section{Không gian trạng thái (State Space)}
Trong trò chơi Bloxorz, không gian trạng thái $S$ bao gồm:
\begin{enumerate}
    \item \textbf{Định hướng và vị trí khối hộp}:
    Khối có kích thước $1 \times 1 \times 2$, tại mỗi bước tồn tại ở một trong ba trạng thái định hướng $O$:
    \begin{itemize}
        \item \textbf{Đứng (Standing - S)}: Khối chiếm 1 ô duy nhất tại tọa độ $(r, c)$.
        \item \textbf{Nằm ngang (Lying Horizontally - H)}: Khối chiếm 2 ô liền kề $(r, c)$ và $(r, c + 1)$.
        \item \textbf{Nằm dọc (Lying Vertically - V)}: Khối chiếm 2 ô liền kề $(r, c)$ và $(r + 1, c)$.
    \end{itemize}
    Trường hợp phân tách khối (Split state): Trạng thái lưu tọa độ của cả hai khối lập phương $1 \times 1 \times 1$: $(r_1, c_1)$ và $(r_2, c_2)$ kèm con trỏ chỉ định khối đang được điều khiển.
    
    \item \textbf{Trạng thái cầu (Bridges Configuration)}:
    Giả sử màn chơi có $k$ cây cầu, trạng thái của các cây cầu được biểu diễn dưới dạng tuple hoặc bitmask nhị phân $B = (b_1, b_2, \dots, b_k)$ với $b_i \in \{0, 1\}$ (0: Đóng / Void, 1: Mở / Sàn).
\end{enumerate}
Do đó, một trạng thái tổng quát được biểu diễn chính xác là một bộ:
\begin{equation}
    s = (O, \text{BlockPosition}, B)
\end{equation}

\section{Không gian hành động (Actions)}
Tại mỗi bước, tập hành động hợp lệ $A = \{\text{Up}, \text{Down}, \text{Left}, \text{Right}\}$. Khi khối bị tách thành 2 khối nhỏ, hành động mở rộng thêm thao tác hoán đổi quyền điều khiển (\texttt{SwitchControl}).

\section{Mô hình chuyển trạng thái (Transition Model)}
Bảng \ref{tab:transition_rules} mô tả quy luật biến đổi tất định vị trí và định hướng khi thực hiện hành động lăn từ tọa độ tham chiếu $(r, c)$:

\begin{table}[htbp]
\centering
\caption{Quy luật chuyển đổi trạng thái khi lăn khối 1×1×2}
\label{tab:transition_rules}
\small
\begin{tabular}{lcccc}
\toprule
\textbf{Trạng thái hiện tại} & \textbf{Up} & \textbf{Down} & \textbf{Left} & \textbf{Right} \\
\midrule
Đứng $S(r, c)$ & $V(r-2, c)$ & $V(r+1, c)$ & $H(r, c-2)$ & $H(r, c+1)$ \\
Nằm ngang $H(r, c)$ & $H(r-1, c)$ & $H(r+1, c)$ & $S(r, c-1)$ & $S(r, c+2)$ \\
Nằm dọc $V(r, c)$ & $S(r-1, c)$ & $S(r+2, c)$ & $V(r, c-1)$ & $V(r, c+1)$ \\
\bottomrule
\end{tabular}
\end{table}

\section{Kiểm tra tính hợp lệ \& Trạng thái kết thúc}
\begin{itemize}
    \item \textbf{Nước đi không hợp lệ (Illegal Move)}: Một hành động lăn bị coi là không hợp lệ nếu bất kỳ phần nào của khối chiếm ô ngoài bản đồ, ô trống (Void), ô cầu đang ở trạng thái Đóng ($b_i = 0$), hoặc khối rơi vào trạng thái Đứng trên ô màu cam dễ vỡ (Fragile Tile).
    \item \textbf{Kiểm tra đích (Goal Test)}: Trạng thái đích đạt được khi và chỉ khi khối ở trạng thái \textbf{Đứng (Standing)} tại tọa độ ô mục tiêu (Goal Hole).
\end{itemize}

\section{Quản lý trạng thái lặp (Cycle Detection \& Visited Set)}
Để đảm bảo thuật toán không bị rơi vào chu trình vô hạn, cấu trúc dữ liệu \texttt{visited} được cài đặt bằng bảng băm (\texttt{set} trong Python) với khóa băm (hashable key) là tuple bất biến:
\begin{equation}
    \text{key} = (O, \text{tuple}(\text{pos}), \text{tuple}(B))
\end{equation}
